\section{General hyperbolic recurrences without residue-ring splitting}

Let $A\in\mathrm{SL}(2,\mathbb Z)$ have $|\operatorname{Tr}A|>2$.
Put $a=|\operatorname{Tr}A|$, $s=\operatorname{sgn}(\operatorname{Tr}A)$,
and define $u_0=0$, $u_1=1$, $u_{j+1}=au_j-u_{j-1}$.
For an odd prime $K>3$ not dividing $a-2s$, set
\[
 m=\frac{K-1}{2},\qquad N_K=u_{m+1}+s u_m,\qquad
 c_K=\left(\frac{a-2s}{K}\right),\qquad
 D_K=\frac{N_K-c_K}{2},\quad n_K=\frac{D_K}{K}.
\]
The parentheses defining $c_K$ denote the Legendre symbol. The exact identity
\[
 A^K-I=s^mN_K A^m(A-I)
\]
implies $A^K=I$ modulo $N_K$. All nontrivial powers have no nonzero fixed
vector: $A-I$ is invertible modulo $N_K$, and Bezout for $j,K$ transfers a
fixed vector of $A^j$ to one of $A$. The modulus is odd; it need not be prime,
squarefree, or split for the characteristic polynomial of $A$.

For exact covariance of integer Fourier symbols at arbitrary $A$, use
\[
 Q_N^*(v)=(-1)^{v_1v_2+v_1+v_2}W_N(v).
\]
This quadratic refinement modulo two is invariant under every invertible
linear transformation of $\mathbb F_2^2$. It is an explicit quantization
characteristic, not an assertion that the earlier characteristic gives the
same finite operator for all $A$.

\paragraph{Finite spectral theorem.}
There is a scalar normalization of the Weyl implementer $U_N$ such that on
$\operatorname{ran}(I-c_KR)/2$ each $K$-th root of unity has multiplicity
exactly $n_K$. On the other parity sector the nontrivial roots have the same
multiplicity, while the root one has multiplicity $n_K+c_K$.
In particular, the uniform sector is even when $c_K=-1$.

The proof does not diagonalize $A$ over $\mathbb Z/N_K\mathbb Z$.
Weyl covariance gives
\[
 |\operatorname{Tr}U_N^j|^2=\#\ker(A^j-I)=1,\quad 0<j<K,
 \qquad |\operatorname{Tr}(RU_N^j)|^2=1.
\]
After normalizing $U_N^K=I$, the second trace is an integral virtual
character of the cyclic group with squared Fourier norm one and value one
at zero. It is therefore a single character. A root-of-unity phase makes
it identically one. If $M_l$ are the full multiplicities, then
\[
 \sum_l M_l=N_K,\qquad
 \sum_l M_l^2=\frac{N_K^2+K-1}{K}.
\]
The least squared sum of integers of a prescribed total forces all
multiplicities to differ by at most one, with exactly one exceptional root.
Parity locates that root at one and gives the asserted sector. The argument
also works for any odd recurrence order whose nontrivial powers are
fixed-point-free, not just a prime order.

For example, trace $+4$ and $K=11$ give $N=989=23\cdot43$ and an even
uniform sector of dimension $495$, with multiplicity $45$. The discriminant
is split modulo $23$ and nonsplit modulo $43$. Trace $-6$ and $K=7$ instead
give $N=169=13^2$ and an odd uniform sector of dimension $84$, with
multiplicity $12$. Its discriminant is nonsplit modulo $13$.

\paragraph{Fixed classical strength.}
Let $S=A^{\mathsf T}$, choose a real even smooth $g_0$ whose support is
disjoint from its $S$-preimage, and put $g=g_0-\int g_0$, $f=g-g\circ S$.
With $\hbar_N=(2\pi N)^{-1}$, define
\[
 V_{\kappa,N}=e^{i\kappa\operatorname{Op}_N^*(f)/\hbar_N}U_N,
 \qquad J_{\kappa,N}=e^{i\kappa\operatorname{Op}_N^*(g)/\hbar_N}.
\]
For every positive odd $j$, the sine Taylor remainder and disjoint derivative
supports give
\[
 \|J_{\kappa,N}^*V_{\kappa,N}J_{\kappa,N}-U_N\|_{\mathrm{op}}
 \leq\frac{4\kappa^2\pi^{j+2}}{j!}
 M_j(g)M_j(g\circ S)N^{2-j},
\]
where $M_j(g)=\sum_v|\widehat g(v)||v|_2^j$. Hence this difference is
$O(N^{-L})$ for every $L$, uniformly on compact classical-strength intervals.
Telescoping preserves that accuracy on every polynomial time window,
including the Heisenberg scale. No smallness of $\kappa N$ is assumed.

The associated classical map is exactly
\[
 S\Phi_f^{-\kappa}=\Phi_g^\kappa S\Phi_g^{-\kappa}.
\]
It is nonaffine and Anosov for nonzero $\kappa$, but retains the arithmetic
structure by smooth symplectic conjugacy. It is therefore not a generic
arithmetic-free BGS example.

\paragraph{Local laws.}
On the uniform parity sector there are $K$ clusters of exactly $n_K$ levels,
with widths smaller than every inverse power of $N_K$. For each fixed
factorial order $p\geq2$ and Schwartz test $\Phi$ on $\mathbb R^{p-1}$,
\[
 R_N^{(p)}(\Phi)=(n_K-1)_{p-1}\Phi(0)+O(N_K^{-L})
\]
for every $L$. At a uniformly random origin, a microscopic interval of
length $x$ is empty with probability $1-x/n_K+O(N_K^{-L})$ and contains
$n_K$ levels with probability $x/n_K+O(N_K^{-L})$. The mean count is exactly
$x$, while its variance is $xn_K-x^2+O(N_K^{-L})$. The local point process
converges to the empty measure; the unit mean intensity is not uniformly
integrable. This is not a Poisson process.

The detailed proofs, including the analytic rigidity boundary and the
quantitative one-mode sublevel estimate for the second resonance, are in the
companion mathematical notes. None of these conclusions is a positive
universality theorem for a general chaotic system or a claim of established
world priority.
